\documentclass[preview]{standalone}

\usepackage{amsmath}
\usepackage{amssymb}
\usepackage{stellar}
\usepackage{definitions}
\usepackage{bettelini}

\begin{document}

\title{Stellar}
\id{linearalgebra-associated-matrices}
\genpage

\section{Matrix-associated maps}

\begin{snippetdefinition}{matrix-associated-linear-map-definition}{Linear map associated to a matrix}
    Let \(\mathbb{F}\) be a \field, and let
    \(A\in\matrices_{m\times n}(\mathbb{F})\). The
    \emph{linear map associated to \(A\)} is the \function
    \[
        L_A\colon \mathbb{F}^n\fromto\mathbb{F}^m,
        \qquad
        L_A(x)=Ax.
    \]
\end{snippetdefinition}

\begin{snippetproposition}{matrix-associated-linear-map-is-linear}{The associated map is linear}
    Let \(\mathbb{F}\) be a \field, let
    \(A\in\matrices_{m\times n}(\mathbb{F})\), and let
    \(L_A\colon \mathbb{F}^n\fromto\mathbb{F}^m\) be defined by
    \(L_A(x)=Ax\). Then
    \[
        L_A\in\setoflineartransformations(\mathbb{F}^n,\mathbb{F}^m).
    \]
\end{snippetproposition}

\begin{snippetproof}{matrix-associated-linear-map-is-linear-proof}{matrix-associated-linear-map-is-linear}{The associated map is linear}
    For every \(x,y\in\mathbb{F}^n\),
    \[
        L_A(x+y)=A(x+y)=Ax+Ay=L_A(x)+L_A(y).
    \]
    For every \(\lambda\in\mathbb{F}\) and every \(x\in\mathbb{F}^n\),
    \[
        L_A(\lambda x)=A(\lambda x)=\lambda Ax=\lambda L_A(x).
    \]
    Thus \(L_A\) is linear.
\end{snippetproof}

\begin{snippetproposition}{matrix-associated-map-composition}{Composition and matrix multiplication}
    Let \(\mathbb{F}\) be a \field, let
    \[
        A\in\matrices_{m\times n}(\mathbb{F}),
        \qquad
        B\in\matrices_{n\times p}(\mathbb{F}).
    \]
    Then
    \[
        L_{AB}=L_A\circ L_B.
    \]
\end{snippetproposition}

\begin{snippetproof}{matrix-associated-map-composition-proof}{matrix-associated-map-composition}{Composition and matrix multiplication}
    Both sides are \function[functions] from \(\mathbb{F}^p\) to
    \(\mathbb{F}^m\). For every \(x\in\mathbb{F}^p\),
    \[
        (L_A\circ L_B)(x)=L_A(Bx)=A(Bx)=(AB)x=L_{AB}(x),
    \]
    where associativity of \matrix multiplication is used.
\end{snippetproof}

\section{Kernel and image}

\begin{snippetproposition}{matrix-associated-kernel-and-image}{Kernel and image of the associated map}
    Let \(\mathbb{F}\) be a \field, let
    \(A\in\matrices_{m\times n}(\mathbb{F})\), and let
    \(L_A\colon \mathbb{F}^n\fromto\mathbb{F}^m\) be defined by
    \(L_A(x)=Ax\). Then
    \[
        \ltkernel L_A
        =
        \{x\in\mathbb{F}^n\suchthat Ax=0\},
    \]
    and
    \[
        \ltimage L_A=\columnspace(A).
    \]
\end{snippetproposition}

\begin{snippetproof}{matrix-associated-kernel-and-image-proof}{matrix-associated-kernel-and-image}{Kernel and image of the associated map}
    By definition,
    \[
        x\in\ltkernel L_A
        \Longleftrightarrow
        L_A(x)=0
        \Longleftrightarrow
        Ax=0.
    \]
    For the image, let \(A^1,\ldots,A^n\) be the columns of \(A\). If
    \(x=(x_1,\ldots,x_n)^t\), then
    \[
        L_A(x)=Ax=x_1A^1+\cdots+x_nA^n.
    \]
    Therefore the values of \(L_A\) are exactly the
    \linearcombination[linear combinations] of the columns of \(A\), so
    \[
        \ltimage L_A=\linearspan(A^1,\ldots,A^n)=\columnspace(A).
    \]
\end{snippetproof}

\begin{snippetcorollary}{matrix-associated-rank-nullity}{Rank-nullity for a matrix-associated map}
    Let \(\mathbb{F}\) be a \field, and let
    \(A\in\matrices_{m\times n}(\mathbb{F})\). Then
    \[
        n=
        \lineardim\{x\in\mathbb{F}^n\suchthat Ax=0\}
        +
        \mrank A.
    \]
\end{snippetcorollary}

\begin{snippetproof}{matrix-associated-rank-nullity-proof}{matrix-associated-rank-nullity}{Rank-nullity for a matrix-associated map}
    Apply rank-nullity to
    \(L_A\colon\mathbb{F}^n\fromto\mathbb{F}^m\). By the previous
    proposition, \(\ltkernel L_A=\{x\in\mathbb{F}^n\suchthat Ax=0\}\), and
    \(\ltrank L_A=\mrank A\).
\end{snippetproof}

\section{Matrices of linear maps}

\begin{snippetdefinition}{linear-map-associated-matrix-definition}{Matrix associated to a linear map}
    Let \(V\) and \(W\) be finite-dimensional \vectorspace[vector spaces] over
    a \field \(\mathbb{F}\). Let
    \[
        \mathcal{B}=\{v_1,\ldots,v_n\}
    \]
    be a \basis of \(V\), let
    \[
        \mathcal{C}=\{w_1,\ldots,w_m\}
    \]
    be a \basis of \(W\), and let
    \(f\in\setoflineartransformations(V,W)\). For each
    \(j\in\{1,\ldots,n\}\), write
    \[
        f(v_j)=a_{1j}w_1+\cdots+a_{mj}w_m.
    \]
    The \emph{matrix associated to \(f\) with respect to
    \(\mathcal{B}\) and \(\mathcal{C}\)} is
    \[
        M(f,\mathcal{B},\mathcal{C})=(a_{ij})
        \in\matrices_{m\times n}(\mathbb{F}).
    \]
    Its \(j\)-th column is the coordinate column
    \[
        [f(v_j)]_{\mathcal{C}}.
    \]
\end{snippetdefinition}

\begin{snippet}{associated-matrix-column-interpretation}
    The columns of \(M(f,\mathcal{B},\mathcal{C})\) are the values of \(f\)
    on the \basis of the domain, expressed in the \basis of the codomain.
    This is the same principle as for a \matrix \(A\): the columns of \(A\)
    are the images of the canonical \basis vectors under \(L_A\).
\end{snippet}

\begin{snippetproposition}{associated-matrix-linearity}{Linearity of the associated matrix construction}
    Let \(V\) and \(W\) be finite-dimensional \vectorspace[vector spaces] over
    a \field \(\mathbb{F}\), with \basis[bases] \(\mathcal{B}\) and
    \(\mathcal{C}\). If
    \[
        f,g\in\setoflineartransformations(V,W),
        \qquad
        \lambda\in\mathbb{F},
    \]
    then
    \[
        M(f+g,\mathcal{B},\mathcal{C})
        =
        M(f,\mathcal{B},\mathcal{C})
        +
        M(g,\mathcal{B},\mathcal{C}),
    \]
    and
    \[
        M(\lambda f,\mathcal{B},\mathcal{C})
        =
        \lambda M(f,\mathcal{B},\mathcal{C}).
    \]
\end{snippetproposition}

\begin{snippetproof}{associated-matrix-linearity-proof}{associated-matrix-linearity}{Linearity of the associated matrix construction}
    Let
    \[
        M(f,\mathcal{B},\mathcal{C})=(a_{ij}),
        \qquad
        M(g,\mathcal{B},\mathcal{C})=(b_{ij}).
    \]
    Then
    \[
        f(v_j)=\sum_{i=1}^m a_{ij}w_i,
        \qquad
        g(v_j)=\sum_{i=1}^m b_{ij}w_i.
    \]
    Hence
    \[
        (f+g)(v_j)
        =
        \sum_{i=1}^m(a_{ij}+b_{ij})w_i,
    \]
    so the \(j\)-th column of \(M(f+g,\mathcal{B},\mathcal{C})\) is the sum
    of the \(j\)-th columns of the two associated \matrix[matrices]. The scalar
    identity follows similarly from
    \[
        (\lambda f)(v_j)=\sum_{i=1}^m(\lambda a_{ij})w_i.
    \]
\end{snippetproof}

\begin{snippettheorem}{linear-maps-matrices-isomorphism}{Linear maps and matrices}
    Let \(V\) and \(W\) be finite-dimensional \vectorspace[vector spaces] over
    a \field \(\mathbb{F}\), with
    \[
        \lineardim V=n,
        \qquad
        \lineardim W=m.
    \]
    Fix \basis[bases] \(\mathcal{B}\) of \(V\) and \(\mathcal{C}\) of \(W\). The map
    \[
        \mathcal{M}_{\mathcal{B},\mathcal{C}}\colon
        \setoflineartransformations(V,W)
        \fromto
        \matrices_{m\times n}(\mathbb{F}),
        \qquad
        f\mapsto M(f,\mathcal{B},\mathcal{C}),
    \]
    is a \linearisomorphism.
\end{snippettheorem}

\begin{snippetproof}{linear-maps-matrices-isomorphism-proof}{linear-maps-matrices-isomorphism}{Linear maps and matrices}
    The previous proposition shows that
    \(\mathcal{M}_{\mathcal{B},\mathcal{C}}\) is linear.

    If \(M(f,\mathcal{B},\mathcal{C})=0\), then \(f(v_j)=0_W\) for every
    \basis vector \(v_j\) of \(V\). Since a \lineartransformation is
    determined by its values on a \basis, \(f=0\). Hence
    \(\mathcal{M}_{\mathcal{B},\mathcal{C}}\) is \injective.

    To prove surjectivity, let
    \(A=(a_{ij})\in\matrices_{m\times n}(\mathbb{F})\). Define values on the
    \basis \(\mathcal{B}\) by
    \[
        f(v_j)=a_{1j}w_1+\cdots+a_{mj}w_m.
    \]
    By linear extension from a \basis, there exists a unique
    \lineartransformation \(f\colon V\fromto W\) with these values. By
    construction,
    \[
        M(f,\mathcal{B},\mathcal{C})=A.
    \]
    Therefore \(\mathcal{M}_{\mathcal{B},\mathcal{C}}\) is \surjective and
    hence a \linearisomorphism.
\end{snippetproof}

\begin{snippetcorollary}{linear-transformations-space-dimension}{Dimension of the space of linear maps}
    Let \(V\) and \(W\) be finite-dimensional \vectorspace[vector spaces] over
    a \field \(\mathbb{F}\), with
    \[
        \lineardim V=n,
        \qquad
        \lineardim W=m.
    \]
    Then
    \[
        \lineardim\setoflineartransformations(V,W)=mn.
    \]
\end{snippetcorollary}

\begin{snippetproof}{linear-transformations-space-dimension-proof}{linear-transformations-space-dimension}{Dimension of the space of linear maps}
    For any choice of \basis[bases] \(\mathcal{B}\) of \(V\) and \(\mathcal{C}\) of
    \(W\),
    \[
        \setoflineartransformations(V,W)
        \simeq
        \matrices_{m\times n}(\mathbb{F}).
    \]
    The latter \vectorspace has dimension \(mn\).
\end{snippetproof}

\section{Composition}

\begin{snippettheorem}{associated-matrix-composition-theorem}{Matrix of a composition}
    Let \(U,V,W\) be finite-dimensional \vectorspace[vector spaces] over a
    \field \(\mathbb{F}\). Let
    \[
        f\in\setoflineartransformations(U,V),
        \qquad
        g\in\setoflineartransformations(V,W).
    \]
    If \(\mathcal{B}\), \(\mathcal{C}\), and \(\mathcal{D}\) are \basis[bases] of
    \(U\), \(V\), and \(W\), respectively, then
    \[
        M(g\circ f,\mathcal{B},\mathcal{D})
        =
        M(g,\mathcal{C},\mathcal{D})\,
        M(f,\mathcal{B},\mathcal{C}).
    \]
\end{snippettheorem}

\begin{snippetproof}{associated-matrix-composition-theorem-proof}{associated-matrix-composition-theorem}{Matrix of a composition}
    Write
    \[
        \mathcal{B}=\{u_1,\ldots,u_n\},
        \qquad
        \mathcal{C}=\{v_1,\ldots,v_h\},
        \qquad
        \mathcal{D}=\{w_1,\ldots,w_m\}.
    \]
    Let
    \[
        M(f,\mathcal{B},\mathcal{C})=(a_{ij}),
        \qquad
        M(g,\mathcal{C},\mathcal{D})=(b_{\ell i}).
    \]
    Then
    \[
        f(u_j)=\sum_{i=1}^{h}a_{ij}v_i,
        \qquad
        g(v_i)=\sum_{\ell=1}^{m}b_{\ell i}w_\ell.
    \]
    Therefore
    \begin{align*}
        (g\circ f)(u_j)
        &=
        g\left(\sum_{i=1}^{h}a_{ij}v_i\right)\\
        &=
        \sum_{\ell=1}^{m}
        \left(\sum_{i=1}^{h}b_{\ell i}a_{ij}\right)w_\ell.
    \end{align*}
    The coefficient of \(w_\ell\) in this expression is the
    \((\ell,j)\)-entry of the product
    \(M(g,\mathcal{C},\mathcal{D})M(f,\mathcal{B},\mathcal{C})\).
\end{snippetproof}

\section{Endomorphisms}

\begin{snippetproposition}{endomorphisms-square-matrices-algebra-isomorphism}{Endomorphisms and square matrices}
    Let \(V\) be a finite-dimensional \vectorspace over a \field
    \(\mathbb{F}\), with \(\lineardim V=n\). If \(\mathcal{B}\) is a
    \basis of \(V\), then
    \[
        \operatorname{End}(V)\fromto\matrices_{n\times n}(\mathbb{F}),
        \qquad
        f\mapsto M(f,\mathcal{B},\mathcal{B}),
    \]
    is an isomorphism of \(\mathbb{F}\)-\algebra[algebras].
\end{snippetproposition}

\begin{snippetproof}{endomorphisms-square-matrices-algebra-isomorphism-proof}{endomorphisms-square-matrices-algebra-isomorphism}{Endomorphisms and square matrices}
    It is a \linearisomorphism by the linear maps and matrices theorem. It
    preserves products because, by the composition formula,
    \[
        M(g\circ f,\mathcal{B},\mathcal{B})
        =
        M(g,\mathcal{B},\mathcal{B})M(f,\mathcal{B},\mathcal{B}).
    \]
    It also sends \(\identityfunc_V\) to the identity \matrix
    \(\identmatrix{n}\). Thus it is an isomorphism of
    \(\mathbb{F}\)-\algebra[algebras].
\end{snippetproof}

\section{Rank}

\begin{snippettheorem}{linear-map-associated-matrix-rank-theorem}{Rank and associated matrix}
    Let \(V\) and \(W\) be finite-dimensional \vectorspace[vector spaces] over
    a \field \(\mathbb{F}\), let
    \(f\in\setoflineartransformations(V,W)\), and choose \basis[bases]
    \(\mathcal{B}\) of \(V\) and \(\mathcal{C}\) of \(W\). Then
    \[
        \ltrank f=\mrank M(f,\mathcal{B},\mathcal{C}).
    \]
\end{snippettheorem}

\begin{snippetproof}{linear-map-associated-matrix-rank-theorem-proof}{linear-map-associated-matrix-rank-theorem}{Rank and associated matrix}
    Let
    \[
        \mathcal{B}=\{v_1,\ldots,v_n\},
        \qquad
        A=M(f,\mathcal{B},\mathcal{C}).
    \]
    Let
    \[
        \Phi_{\mathcal{C}}\colon W\fromto\mathbb{F}^m
    \]
    be the coordinate \linearisomorphism with respect to the \basis
    \(\mathcal{C}\). By definition of the associated \matrix, the \(j\)-th
    column \(A^j\) of \(A\) is
    \[
        A^j=\Phi_{\mathcal{C}}(f(v_j)).
    \]
    Since
    \[
        \ltimage f=\linearspan(f(v_1),\ldots,f(v_n)),
    \]
    applying the isomorphism \(\Phi_{\mathcal{C}}\) gives
    \[
        \Phi_{\mathcal{C}}(\ltimage f)=\linearspan(A^1,\ldots,A^n)
        =
        \columnspace(A).
    \]
    Isomorphisms preserve dimension, so
    \[
        \ltrank f=\lineardim\ltimage f=\lineardim\columnspace(A)=\mrank A.
    \]
\end{snippetproof}

\section{Dimensional consequences}

\begin{snippetexample}{more-equations-than-unknowns-map-example}{More equations than unknowns}
    Let \(A\in\matrices_{3\times2}(\realnumbers)\). Then
    \[
        L_A\colon\realnumbers^2\fromto\realnumbers^3.
    \]
    This \lineartransformation cannot be \surjective. Indeed, rank-nullity
    gives
    \[
        \lineardim\ltimage L_A\leq 2,
    \]
    while \(\lineardim\realnumbers^3=3\). Hence there exists
    \(b\in\realnumbers^3\) such that the system
    \[
        Ax=b
    \]
    has no solution.
\end{snippetexample}

\begin{snippetexample}{more-unknowns-than-equations-map-example}{More unknowns than equations}
    Let \(A\in\matrices_{2\times3}(\realnumbers)\). Then
    \[
        L_A\colon\realnumbers^3\fromto\realnumbers^2.
    \]
    This \lineartransformation cannot be \injective. Indeed,
    \[
        3=\lineardim\ltkernel L_A+\lineardim\ltimage L_A,
    \]
    and
    \[
        \lineardim\ltimage L_A\leq 2.
    \]
    Hence \(\lineardim\ltkernel L_A\geq 1\), so
    \[
        \ltkernel L_A\neq\{0\}.
    \]
    In terms of linear systems, the homogeneous system
    \[
        Ax=0
    \]
    has nonzero solutions.
\end{snippetexample}

\begin{snippetexample}{injective-not-surjective-dimensional-example}{Injective but not surjective by dimension}
    Consider
    \[
        f\colon\realnumbers^2\fromto\realnumbers^3,
        \qquad
        f(x,y)=(3x+2y,\ x-y,\ 2x+7y).
    \]
    The kernel is given by the system
    \[
        \begin{cases}
            3x+2y=0,\\
            x-y=0,\\
            2x+7y=0.
        \end{cases}
    \]
    From \(x-y=0\), we get \(x=y\). Substituting into
    \(3x+2y=0\) gives \(5x=0\), so \(x=y=0\). Hence \(f\) is
    \injective. However,
    \[
        \ltimage f
        =
        \linearspan\bigl((3,1,2),(2,-1,7)\bigr),
    \]
    so \(\lineardim\ltimage f\leq2\), and \(f\) is not \surjective onto
    \(\realnumbers^3\).
\end{snippetexample}

\end{document}
